% Knowledge Graph Constructor: Design and Implementation (about 17 pages). Written in Stage 5.
% Facts about the code cite the final code at R230 (src/kgbuilder/...); numbers have rows in evidence_registry.csv.
\chapter{Knowledge Graph Constructor: Design and Implementation}\label{ch04}

This chapter describes the builder: the part of the system that turns a folder of tables and documents into a
graph. It first states the design principles (\cref{sec:ch04-principles}) and the stages
(\cref{sec:ch04-architecture}), then follows the data through the structured and text paths
(\cref{sec:ch04-structured,sec:ch04-text,sec:ch04-extraction}). \Cref{sec:ch04-evolution} explains, with the
measurements that motivated it, why the final graph is an \emph{anchor graph} --- an index of identity and
evidence --- rather than a graph of extracted facts. The remaining sections describe that graph: its schema,
how mentions receive their identity, how text is linked to records, how provenance is kept, and how the graph
is navigated. \Cref{ch05} evaluates it.

\section{Objectives and Design Principles}\label{sec:ch04-principles}

The builder has one objective: from mixed tables and text, and a goal stated in one sentence, build a graph in
which the things the data talks about can be found and from which the evidence about each of them can be
reached --- correctly, and without code written for the dataset at hand. Five principles follow from the
requirements of \cref{sec:ch03-requirements}.

\begin{enumerate}[nosep]
\item \textbf{The language model proposes; code decides.} A model is used where something must be proposed
that rules cannot compute: the mapping of tables to a graph, the kinds of facts in the text, the facts and the
mentions themselves, and the identity of near misses. Every proposal is parsed into a typed model and checked by
code before it is used.
\item \textbf{Evidence before storage.} Every item made from text carries a quotation that code finds in its
source chunk; an item whose quotation is not found is rejected, not repaired.
\item \textbf{Records are exact.} A row of a table becomes a node identified by its source key. Text may link
to a record; it never changes or merges one.
\item \textbf{Identity is an edge.} The decision that a mention refers to a record, to a person or to a concept
is stored as an edge carrying its reason, not as a merged node, so it can be inspected and undone.
\item \textbf{Domain neutrality.} Nothing in the code or the prompts names a domain. The examples in the
prompts are invented and checked by tests against the evaluation corpora (\cref{sec:ch04-engineering}).
\end{enumerate}

\section{Architecture Overview}\label{sec:ch04-architecture}

The builder is a pipeline of stages, each a command of the \texttt{kg} command-line tool and each recorded as
one MLflow run. \Cref{tab:builder-stages} lists them in their order; \cref{fig:system-overview} shows where
they sit in the whole system.

\begin{table}[htbp]
\centering
\caption{The builder's stages, in order. ``LLM'' names what a language model proposes; everything else is
deterministic code.}
\label{tab:builder-stages}
\small
\begin{tabularx}{\linewidth}{@{}lXX@{}}
\toprule
Stage & What it does & LLM role and check \\
\midrule
\texttt{profile} & stages the tables and profiles every column in DuckDB & --- \\
\texttt{plan} & maps tables to node labels, keys and relationships & proposes the plan; code validates it, then a critic reviews it \\
\texttt{build} & imports the rows as nodes and relationships & --- (re-validates the plan) \\
\texttt{ingest-text} & reads documents, cuts chunks, embeds them & embeddings only \\
\texttt{text-schema} & names the kinds of things and facts in the text & proposes the schema; code validates, a critic reviews \\
\texttt{extract} & extracts claims with their quotations & proposes claims; code verifies each quotation and argument \\
\texttt{link} & links documents and sections to records; derives implied claims & --- \\
\texttt{mention-pass} & finds every mention of a thing in every chunk & proposes mentions; code filters them \\
\texttt{resolve} & decides what each mention refers to & only for near-miss records and uncertain pairs; code verifies quotations \\
\texttt{attach} & attaches claims and unattached documents to records & --- \\
\bottomrule
\end{tabularx}
\end{table}

\paragraph{Ports and the composition root.} The language model, the embedder and the experiment tracker are
reached through small interfaces (Python \texttt{Protocol} classes), and the concrete objects --- the Gemini
client, the Neo4j driver, the MLflow tracker and the settings --- are built in one place, the command-line entry
point, and passed down. This lets every stage run in tests against fakes, without a network or a paid model.
The Neo4j driver is the one dependency injected directly rather than behind an interface of its own.

\paragraph{Models.} All final builds use Gemini~3.8~Flash for every proposing role, with medium thinking for the
plan and the text schema and low thinking for extraction, the mention pass and identity decisions; embeddings
come from \texttt{gemini-embedding-001}. Two other models were tried and rejected. A local 7-billion-parameter
model could not produce a valid plan, and about 80\,\% of its facts failed the quotation check. DeepSeek~V4.1
Flash built three graphs in one arm of the project, but its text schema on the NHTSA data had no relation from a
part to a vehicle, and claim recall fell from 56/68 with Gemini to 42/68.

\section{Structured Ingestion}\label{sec:ch04-structured}

\paragraph{Staging and profiling.} CSV files are copied as they are; JSON, JSON-lines and NDJSON files are
flattened to CSV; tables of a PostgreSQL database are exported through the same interface, with the keys the
database declares. DuckDB then profiles every column: its type, its number of empty and distinct values, five
sample values, whether it is unique, and --- for text columns --- whether it holds prose (on average at least
120 characters and at least half of its values with more than one sentence). Foreign-key candidates are found
by inclusion: a column whose distinct values are at least 95\,\% contained in a unique column of the same type
in another file. The model never sees the rows themselves, only this profile.

\paragraph{The plan.} The model proposes a \emph{construction plan}: a list of node rules (a source file, a
label, the key column or columns, the properties, an optional name column) and relationship rules (a source file,
a type, and the columns that join a source label to a target label). The plan is a declarative description of
the graph, not code: it is validated, stored as a file a person may review or edit, and executed by a generic
importer. Code checks, among other things, that:
\begin{itemize}[nosep]
\item every column named exists and every key is unique in its file;
\item a composite key (several columns that together identify a row, added for datasets whose rows have no
single key) is complete and unique over all staged rows, and no part of it is unique alone;
\item every relationship joins defined labels through existing columns;
\item a name column is not a code-like identifier;
\item files that the foreign-key analysis joins end up in one connected part of the plan.
\end{itemize}

\paragraph{The refinement loop.} The plan --- and later the text schema --- is produced by a three-round loop
(\cref{alg:refine}). The critic is asked only when code has found nothing wrong, and is told what code already
guarantees, so that it judges the modelling rather than repeating the checks.

\begin{algorithm}[htbp]
\caption{Propose, validate, critique (\texttt{llm/refine.py}).}\label{alg:refine}
\small
\begin{enumerate}[nosep,label=\arabic*.]
\item feedback $\gets$ none
\item \textbf{for} round $= 1, 2, 3$:
\item \quad proposal $\gets$ model(profile, goal, feedback)
\item \quad issues $\gets$ validate(proposal) \hfill \emph{code}
\item \quad \textbf{if} no issues: issues $\gets$ critic(proposal) \hfill \emph{model, asked to say ``valid'' or ``retry''}
\item \quad \textbf{if} no issues: \textbf{return} proposal (accepted)
\item \quad feedback $\gets$ the proposal and its issues
\item \textbf{return} the last proposal, rejected (written for review; the stage fails)
\end{enumerate}
\end{algorithm}

\paragraph{Import.} For each node rule the importer creates a uniqueness constraint on the key and writes the
rows in batches of 1{,}000 with \texttt{MERGE}, so that re-running the import never duplicates a node. Values are
typed by DuckDB; empty cells are left out rather than written as nulls. Relationships are matched by their key
columns and also written with \texttt{MERGE}. Prose columns are turned into \emph{record documents}: one document
per record, with one section per prose cell, which then follow the text path and are linked to their record by
its key.

\section{Text Ingestion}\label{sec:ch04-text}

Documents in Markdown, plain text and PDF are read and cut into chunks. A chunk boundary is first placed at every
section break; a section longer than 1{,}500 characters is cut at paragraph boundaries, then at whitespace; pieces
shorter than 200 characters are joined to the next one. Chunks do not overlap. Each chunk keeps its document's
first heading as \emph{context}, which later checks may use to resolve a name the chunk itself leaves implicit.
Chunks are written to Neo4j as \texttt{:Chunk} nodes linked to their \texttt{:Document} by \texttt{PART\_OF} and
to each other in reading order by \texttt{NEXT\_CHUNK}, and their embeddings are stored in a Neo4j vector index
(cosine similarity). Full-text indexes over chunks and claims are added by a separate indexing stage, used by
the query engine.

\section{Text Schema and Claim Extraction}\label{sec:ch04-extraction}

\paragraph{The text schema.} From the chunks (all of them up to 200{,}000 characters, else an even sample), the
plan's labels and the goal, the model proposes the \emph{entity types} and \emph{fact types} of the text. Each fact
type has a predicate, a subject type and an object type; each entity type has an \emph{identity class}:
\emph{keyed} (it is a record of the plan, named by its labels and key attributes), \emph{individual} (one
particular thing named only in text) or \emph{concept} (a kind, a part, a state). A built-in \texttt{Value} type
holds numbers. Code checks the names, that every referenced type exists and is used, that a keyed type points to
real plan labels and columns, and that a type named like a plan label is keyed; the critic then looks for
overlapping types, claims no fact type could hold, and types the goal does not need. The prompt asks for at most
12 entity types and 20 fact types.

\paragraph{Extraction.} For each chunk the model returns claims: a subject, a predicate, an object, their types,
a verbatim quotation, and assertion fields --- polarity, truth (affirmed or negated) with the negation words,
modality (actual, possible or conditional) with the hedge or condition words, and a time if one is written. Code
then verifies every claim and rejects it, with a recorded reason, if
\begin{itemize}[nosep]
\item its predicate is not an extractable fact type of the schema;
\item its quotation, after normalisation (case, accents, Markdown marks and spacing), is not a substring of the
chunk;
\item a subject or object is empty, a pronoun, the same as the other, or not written in the chunk or its context;
\item a value is not a number found in the quotation, or a time is not in it;
\item a negation, hedge or condition word is not a whole word of the quotation.
\end{itemize}
An optional second pass per chunk (\emph{gleaning}) sees the claims already accepted and asks for those that
were missed, under the same checks. The final builds of the furniture, NHTSA and generality datasets used two
passes; the NTSB build used the default single pass. Each accepted claim is stored as an
\texttt{:Observation} node linked to the mentions of its subject and object and, by \texttt{FROM}, to its chunk.

\section{Design Evolution: From Exhaustive Facts to an Anchor Graph}\label{sec:ch04-evolution}

The graph described in the rest of this chapter is the fourth of four representations built during the project.
Each change was prompted by a measured failure of the previous one. This section presents that sequence as the
argument for the final design; \cref{sec:ch05-representations} gives the full measurements.

\paragraph{Representation 1: claims as edges.} The first graph stored every extracted claim as an edge between
nodes shared by the whole corpus --- for example, \emph{drawer rails} $\xrightarrow{\textsc{has\_defect}}$
\emph{rough, uneven surfaces}. Extraction itself was precise: on the furniture reviews, 170 of 170 stored claims
were judged right, and 89 of the 96 reference claims were found; on the NHTSA complaints, 135 of 146 and 56 of 68
(\cref{tab:representations}). But a path through such a graph is read as a fact, and paths crossed from one
thing's evidence into another's. An audit of the paths a question would follow found 36 of 100 false on the
furniture data and 14 of 56 on the NHTSA data. \Cref{fig:leak}~(A) shows the typical case: the defect stated in
one product's review, stored on the shared part node, is reached from a different product whose own review
praises the part. Moreover, on the NHTSA data the text did not reach the records at all: no extracted entity was
linked to a vehicle or a recall.

\begin{sidewaysfigure}
\centering
\includegraphics[width=\linewidth]{diagrams/construction/r63_r230/r230_diag1_leak_and_witness_rule.pdf}
\caption{The leak of a claim stored on a shared node (A) and the witness rule of the anchor graph (B), on the
furniture data. The words are those of the reviews.}
\label{fig:leak}
\end{sidewaysfigure}

\paragraph{Representation 2: claims as nodes.} Making each claim a node of its own, attached to the thing it is
about, removed the leak: every audited path was true (89 of 89 and 107 of 107). It exposed a second limit
instead. Measured against what a careful reader finds in 40 random sentences, the stored claims covered only
21 of 67 statements on the NHTSA data and 35 of 72 on the furniture data. Most misses had no fact type in the
schema that could hold them --- a repair, an inspection finding --- so they could not be extracted at all.

\paragraph{Representation 3: a layered model.} The third arm kept claims as nodes and added layers: mentions
linked to canonical entities, attachment of claims to records, assertion fields (truth, modality, condition), and
a query layer of fixed primitives over them. Each layer was measured, and each worked by its own measure: for
example, negated claims were stored exactly in 13 of 15 cases on the furniture data, up from 10 of 15. But the
answers to questions barely moved. A change to the reading of claims changed no answer at all. Compared, on the
same questions, with a system that queries only the records and retrieves passages by vector similarity, the claim
graph won on the NHTSA data (60 against 50 of 68 questions, McNemar 13/3, $p = 0.021$), tied on the furniture
data (34 against 34 of 68), and lost on the generality corpus (17 against 25 of 41, 0/8, $p = 0.008$), where plain
vector retrieval did best (27 of 41). An analysis of the 69 wrong answers of the layered system attributed 10 to
facts no schema type could hold and 41 to the query layer, not to the graph.

\paragraph{What the measurements said.} Four conclusions were drawn on 6~October~2026:
\begin{enumerate}[nosep]
\item \emph{The records do the exact work.} Where the graph won, it won through record questions; no question
type separated it from records plus vector retrieval on any dataset.
\item \emph{A claim schema limits generality.} Every new kind of statement needs a fact type, and a corpus of
another domain brings statements no proposed schema anticipates.
\item \emph{Claims cannot be counted exactly anyway.} With a third to a half of the text's statements stored, an
answer that counts over text must read the text itself.
\item \emph{What was worth keeping is identity, linking and provenance:} mentions linked to entities by a
reversible edge, text linked to records, and a verbatim quotation for every item made from text --- which held
in every audit.
\end{enumerate}

\paragraph{Representation 4: the anchor graph.} The fourth design therefore changed the graph's job. Its goal, as
fixed in the design document, is that \emph{the graph lets a reader find the entities, concepts and records a need
is about, and leads from each of them to every chunk, document and record that concerns it, and to nothing that
does not.} The graph stores who and what is mentioned where, the records, and the links between them; it does
not rely on what the text claims. Understanding what a chunk says is left to the reader of the chunk --- in the
final system, the query agent of \cref{ch06}. One rule makes the change concrete:

\begin{quote}
\emph{A shared node helps to find chunks; it never proves a fact.} Two things are connected only by a record
relation or by a chunk that mentions both.
\end{quote}

In \cref{fig:leak}~(B), the question about defective drawer rails starts at the shared concept, which leads to the
chunks that mention it; each chunk leads to the product it is about; and the bed's chunk is rejected by reading
it. A wrong merge of two concepts now costs one extra chunk to read, not a false fact.

\paragraph{Kept, frozen and dropped.} The anchor graph reused what the earlier arms had built and measured:
chunking and embeddings, the plan and the record graph, the linking of text to records, mentions with
\texttt{REFERS\_TO} identity edges, document attachment and the provenance checks. The claim layer was
\emph{frozen}: it is still extracted, stored and evaluated (\cref{sec:ch05-extraction}), but it is not used to
navigate, and an ablation compares navigation with and without it (\cref{sec:ch05-representations}). Building
more claim meaning into the graph --- events, roles, speakers, open predicates --- was dropped.

\paragraph{The evaluation changed with the goal.} Since the graph is now an index, it is judged as an index: can a
target be found, does it lead to the right evidence and only to that, are connected things connected, and is every
link from text to records right. These are the criteria C0--C9 of \cref{sec:ch05-method}, fixed in the same
document before they were measured.

\section{The Anchor Graph Schema}\label{sec:ch04-schema}

\Cref{fig:anchor-layers} shows the final graph. It has four layers and one attachment.

\begin{sidewaysfigure}
\centering
\includegraphics[width=\linewidth]{diagrams/construction/r63_r230/r230_diag2_anchor_graph_layers.pdf}
\caption{The anchor graph's node labels, edges and walks. Records come from the tables under the proposed plan;
every other node is made from text.}
\label{fig:anchor-layers}
\end{sidewaysfigure}

\begin{description}[nosep]
\item[Lexical layer.] \texttt{(:Document)<-[:PART\_OF]-(:Chunk)-[:NEXT\_CHUNK]->(:Chunk)}.
\item[Mention layer.] \texttt{(:Chunk)-[:MENTIONS]->(:Mention)}. A mention is one name of one type in one
document, with the class the text states for it (a particular thing or a kind).
\item[Identity layer.] \texttt{(:Mention)-[:REFERS\_TO]->(\emph{record} | :Individual | :Concept)}, each edge with
its reason, score, evidence and the deciding party (code or a model).
\item[Record layer.] Records under the labels of the plan, joined by the plan's typed relationships.
\item[Attachment.] \texttt{(:Document | :Chunk)-[:ABOUT]->(\emph{record})}: a document or a section is about one
record.
\end{description}

Three kinds of thing are told apart, because their identity works differently (\cref{tab:identity-classes}).

\begin{table}[htbp]
\centering
\caption{The three classes of thing and their identity.}
\label{tab:identity-classes}
\small
\begin{tabularx}{\linewidth}{@{}lXX@{}}
\toprule
Class & Identity & Who may join it with another \\
\midrule
Record & its source key & nobody; text only links to it \\
Individual & its type, name and document & the resolver, across documents, only with evidence \\
Concept & its type and normalised name & the resolver, under guards against wrong merges \\
\bottomrule
\end{tabularx}
\end{table}

\section{Mentions and Identity}\label{sec:ch04-identity}

\paragraph{The mention pass.} The claims of \cref{sec:ch04-extraction} name only the things that take part in a
claim of a known type. The anchor graph needs every thing a question might start from, so a separate pass asks
the model, per chunk, for ``the entries an index of the text would have'': named things, kinds the text says
something about, faults written as verbs, and things the text says are absent --- but not describing words alone,
generic words, everyday acts, the document or its writer. For each it gives the name, whether it is a particular
thing or a kind, and a type of the schema. Code then refuses a finding whose name is not written as whole words in
the chunk, is a quantity, a date, a pronoun or a function word, has more than six words, contains clause
punctuation, or repeats one already listed; and it files the mention under a type only if that type's identity
class can hold the stated class, otherwise under a fallback type for particulars or kinds.

\paragraph{Resolution.} The \texttt{resolve} stage then decides, for every mention, what it refers to
(\cref{fig:identity}). The governing rule is that code decides wherever a rule can, a model is asked only where
rules cannot decide, and a model's answer counts only if code can verify the quotations it gives.

\begin{sidewaysfigure}
\centering
\includegraphics[width=\linewidth]{diagrams/construction/r63_r230/r230_diag3_identity_resolution.pdf}
\caption{Identity resolution per class of thing. Thresholds are the code's settings; spelling scores are
token-sort similarity ratios from 0 to 100.}
\label{fig:identity}
\end{sidewaysfigure}

\paragraph{Mentions to records.} A mention of a keyed type is linked to a record by the first of these tiers that
gives exactly one record:
\begin{enumerate}[nosep]
\item \emph{key}: the record's key is written in the name, and the key may decide alone --- always if it contains
a digit, otherwise only if every other word of the name is a number or one of the record's own attribute values;
\item \emph{name}: the name matches the record's name within the \emph{scope} of the document --- the records
within two relationship hops of what the document is about --- with the same words in any order, or a similarity of
at least 90 and the same words up to their endings;
\item \emph{key in sentence}: exactly one key is written right next to the name in a sentence;
\item \emph{attribute}: among tied records, exactly one has a key attribute value written in a sentence naming
the mention;
\item \emph{variant attribute}: the same, among records whose names are variants of the mention (initials, titles,
short forms).
\end{enumerate}
Mentions no tier settles but for which near misses exist --- names sharing a word stem, similar spellings, a key
in the name that may not decide alone --- go to a \emph{chooser}: the model sees up to five sentences naming the
mention and each candidate's fields and relations, and returns a record and a quotation. Code links only if the
record was among those listed, the quotation is found in one of the mention's chunks and contains the mention's
name, and no other candidate has the same name. A mention the text states as a kind, with no record found, becomes
a concept.

\paragraph{Individuals.} Mentions of individual types are first grouped by the record they link to, if any.
Pairs of the remaining groups are nominated by name variants, by spelling (similarity at least 80), or by
meaning (mutual nearest neighbours of their name embeddings). The model reads up to four sentences naming each
side, with their neighbours, and answers \emph{same}, \emph{different} or \emph{unsure}, quoting one sentence for
each side. Code joins a pair only if both quotations are found on their own side; two different records are never
joined.

\paragraph{Concepts.} Mentions of concept types form one concept per type and normalised name. Two concepts of
the same type are candidates if their spelling similarity is at least 80, or if they are mutual nearest neighbours
($k = 2$) by embedding. Four guards then block a merge: the two have claims of opposite polarity; one sentence
names both as different things; a part-of claim joins them; or one name is the other plus one final word.
Candidates scoring at least 92 are merged by code, the rest are put to the model. Concept merges, unlike the other
two classes, carry no quotation check; the anchor graph tolerates their errors better, since a wrong concept
merge only adds chunks to read (\cref{sec:ch04-evolution}).

\paragraph{Reversibility.} Every outcome is written as a \texttt{REFERS\_TO} edge whose reason names the tier or
the decision (\texttt{key}, \texttt{name}, \texttt{chosen}, \texttt{adjudicated}, \dots) and whose \texttt{by}
property names code or the model. The whole identity layer can be removed by one command, leaving the mentions,
claims and records untouched.

\section{Record Linking and Attachment}\label{sec:ch04-linking}

Identity links \emph{mentions} to records. Two further steps link \emph{documents and chunks} to records, which
define the scopes used above and give the walks of \cref{sec:ch04-navigation} their starting points.

\paragraph{Document and section links.} A document is \texttt{ABOUT} a record when the record's name, of at least
four characters, is written in the document's title; the longest match wins. A record document is about the record
it was made from. A chunk is \texttt{ABOUT} a record when the record's key appears as a whole word in a heading of
the chunk. For a document left without a subject, the attach stage adds a link to a \emph{dominant} record: one
named in at least two sentences and at least twice as often as any other.

\paragraph{Derived claims.} Some fact types are \emph{derived}: their claims are implied by the text's position
rather than written in it (a review's defect belongs to the product the review is about). For these, code writes a
claim from a mention in a chunk to the record its document is about, with the first sentence that names the
subject as its quotation, and only when the record's label fits the fact type's object type.

\paragraph{Claim attachment.} Finally each claim is attached to the records and individuals it concerns, by the
most specific of four routes: a name of the thing written in the claim's quotation sentence; a part-of claim
with the same subject in the same document; the chunk's section; or the document. Of the routes that find
something of one kind, only the most specific is kept. Concepts never hold claims.

\section{Provenance Model}\label{sec:ch04-provenance}

Every element made from text records where it came from (\cref{tab:provenance}). A separate check re-verifies the
finished graph: every claim's quotation is in its chunk, and every link with evidence carries it. In the terms of
the W3C PROV model~\cite{moreau2013provdm}, each claim, mention and link is an entity generated by a recorded
activity (a stage run, logged in MLflow), derived from a chunk or a record, and attributed to code or a named
model.

\begin{table}[htbp]
\centering
\caption{Provenance stored with each element made from text.}
\label{tab:provenance}
\small
\begin{tabularx}{\linewidth}{@{}lX@{}}
\toprule
Element & Provenance \\
\midrule
Claim (\texttt{:Observation}) & its chunk (\texttt{FROM} edge and id), the verbatim quotation, the subject and object names as written, the extractor (model id, or \emph{derived}) \\
Mention & its document; the chunks that mention it (\texttt{MENTIONS}); its stated class \\
\texttt{REFERS\_TO} & reason, score, evidence, and who decided (code or the model id) \\
\texttt{ABOUT} & the matched name; for text-made links also the route and evidence \\
Claim attachment & the route and its evidence (a name, a claim's quotation, a chunk or a document id) \\
Record document & the label, key property and key of its record \\
\bottomrule
\end{tabularx}
\end{table}

\section{The Navigation Contract}\label{sec:ch04-navigation}

The anchor graph promises five walks, the interface between the graph and whatever reads it
(\cref{tab:walks}). A question that needs several hops is a composition of them --- for example, from a complaint's
component to its chunks (W1, W2), from a chunk to its vehicle (W3), and from the vehicle to its recalls (W4) ---
and every hop is a record relation or a chunk, as the witness rule requires.

\begin{table}[htbp]
\centering
\caption{The walks of the navigation contract (\texttt{anchor/navigation.py}).}
\label{tab:walks}
\small
\begin{tabularx}{\linewidth}{@{}llX@{}}
\toprule
Walk & From $\rightarrow$ to & Through \\
\midrule
W1 find & a name or alias $\rightarrow$ its node & exact and normalised names of records, individuals and concepts, then shared words; no model \\
W2 chunks of a node & a node $\rightarrow$ the chunks that concern it & \texttt{REFERS\_TO} and \texttt{MENTIONS} back to chunks; for records and individuals also \texttt{ABOUT} \\
W3 about / names & a chunk $\rightarrow$ the records and individuals it is about or names & \texttt{ABOUT}, else its document's; and its mentions' \texttt{REFERS\_TO}, never to a concept \\
W4 related & a record $\rightarrow$ related records & the plan's typed relationships \\
W5 context & a chunk $\rightarrow$ its neighbours and document & \texttt{NEXT\_CHUNK}, \texttt{PART\_OF} \\
\bottomrule
\end{tabularx}
\end{table}

W3 never follows a concept: a kind is shared by every chunk that names it, so following it from a chunk would join
unrelated chunks through the kind --- exactly the leak of \cref{fig:leak}. A concept remains a valid start (W1, W2).
The walks are implemented once, over an offline snapshot of the graph rebuilt from the build's own records, so
that the evaluation of \cref{ch05} needs neither a database nor a model.

\section{Engineering}\label{sec:ch04-engineering}

\paragraph{Idempotence and safe Cypher.} Every write is a \texttt{MERGE}; derived layers (attachment, identity,
\texttt{ABOUT}) are deleted and recomputed on each run, so a stage can be repeated without leaving stale state.
All values travel as query parameters; labels, relationship types and property names, which cannot be parameters,
pass through a single escaping function.

\paragraph{Caching and retries.} Every model call is cached on a hash of the model, temperature, prompt, response
schema and thinking level. A repeated build is therefore free and gives the same graph, which is how every graph
in this thesis was reloaded for each question set. Failed calls are retried up to three times with growing pauses.

\paragraph{Tracking.} Each stage is an MLflow run nested under a pipeline run, logging its parameters --- including
a hash of each prompt template, so a changed prompt is visible --- its counts, tokens and cost, and its output files.
Tracking failures fall back to a null tracker and never stop a build.

\paragraph{Prompts.} The builder sends nine prompts and two prompt fragments to the model: plan and plan critic,
text schema and schema critic, extraction and its gleaning suffix, the mention pass, the record chooser, the
individual adjudicator and the concept adjudicator. Their examples use invented domains, and tests reject any
prompt that contains a four-word sequence from the evaluation corpora or a domain word of the datasets
(\cref{app:prompt-templates}).

\paragraph{Tests.} The builder's stages, checks and walks are covered by the project's automated tests, which run
against fakes of the model, the embedder and the tracker; tests that need Neo4j run against a separate test
database.
